\section{Introduction}
\label{op:introduction}

A group algebra $K[G]$ consists of finite sums $\sum_g c_g g$, with
$c_g\in K$, multiplied by the group law and bilinearity. A unital ring is
\emph{directly finite} when $ab=1$ implies $ba=1$; it is \emph{stably
finite} when this implication holds in every finite square matrix ring.
Kaplansky's direct-finiteness conjecture asks whether $K[G]$ is directly
finite for every field $K$ and every group $G$, including groups with
torsion. The result here is a negative answer in one specified
odd characteristic.

\begin{theorem}\label{op:main}
Put
\[
 m=\binom{1200}{600},\qquad
 p=\min\{h\in\mathbb Z:h>1,\ h\mid(m!)^2+1\}.
\]
Then $p$ is an odd prime. There are a specified field $K$ of order $p^4$,
a specified finitely generated group $G$ containing torsion, and specified
finite sums $a,b\in K[G]$ satisfying $ab=1$ and $ba\ne1$.
\end{theorem}

The word ``specified'' refers to finite formulas and choices in ordered
finite sets. The field is a quadratic tower, the group is built from
finitely many finite groups and explicit identifications, and the sums
are defined without solving the word problem of that group. The
construction is not claimed to have been evaluated as a practical
computation. It concerns the displayed prime, not every odd prime.

\paragraph{The finiteness problem.}
For finite $G$, the algebra and all its finite matrix rings are
finite-dimensional over $K$. A one-sided inverse between their left
multiplication operators is therefore two-sided. This argument works even
when the characteristic divides $|G|$. The difficulty for an infinite
group is to obtain a corresponding conclusion from finite support alone.

Kaplansky proved stable finiteness in characteristic zero using finite
von Neumann algebras, and asked whether the one-sided-inverse property
persists in positive characteristic~\cite[pp.~122--123]{Kaplansky1972}.
Ara, O'Meara, and Perera proved stable finiteness for free-by-amenable
groups over arbitrary division rings~\cite{AraOMearaPerera2002}.
Elek and Szab\'o~\cite[Corollary~4.7]{ElekSzabo2004} established the
same conclusion for all sofic groups. Such groups admit finite
permutation models that approximate multiplication while keeping
distinct elements separated.
These theorems are constraints on the example, not hypotheses about its
group. In particular, Elek--Szab\'o's theorem implies, after the scalar
construction is completed, that $G$ cannot be sofic.

A separate characteristic-two construction uses incidence structures,
polynomial pairings, a planar group-embedding argument, and character
packing~\cite[Theorem~1.1]{OpenAICharacteristicTwo2026}.
The construction here has a different local mechanism. A high-degree
subspace of a modular translation representation is annihilated on one
side; a weighted pairing gives a compatible annihilation on the other.
The characteristic-two theorem is not an input to this odd-prime
argument, and the two methods retain distinct coefficient hypotheses.

\paragraph{The two quantities that must be compared.}
A matrix $F$ \emph{fits} a simple graph when $F_{ii}=1$ and $F_{ij}=0$
for distinct nonadjacent vertices. We use a triangle-free graph on $m$
vertices and a specified factorization
\[
 F_{ij}=\sum_{T\in\mathcal T}\alpha_i(T)\beta_j(T),\qquad
 |\mathcal T|=t,\qquad \frac{t}{m}<\frac1{100}.
\]
This is an explicit threshold-Kneser instance of the polynomial
constructions of Haviv and Golovnev--Haviv
\cite[Proposition~13]{Haviv2018Theta}
\cite[Theorem~3.2 and Lemma~3.5]{GolovnevHaviv2022}, whose polynomial
ancestry includes Alon--Babai--Suzuki~\cite{AlonBabaiSuzuki1991}.
Low-rank fitting matrices also give storage-code recovery equations.
The triangle-free constructions of Barg--Z\'emor, Huang--Xiang, and
Barg--Schwartz--Yohananov provide related context
\cite{BargZemor2022,HuangXiang2023,BargSchwartzYohananov2024}.
Their results are not used to supply a group-algebra idempotent.

At every vertex we use a copy of a finite group $H$ and an idempotent
$e\in K[H]$. Write $w=|H|$ and $d=\dim_K eK[H]$, the rank of left
multiplication by $e$ on the regular representation. The second quantity
is
\[
 \frac d w>\frac1{96}.
\]
The idempotents satisfy a directed condition: after identifying certain
cyclic subgroups at adjacent vertices, $e_i e_j=0$ for every edge with
$i<j$. There is no assertion about the reverse product. The asymmetry is
what allows a triangular calculation to combine the two strict bounds.

The characteristic-dependent construction appears in
Section~\ref{op:modular-section}. High augmentation degree gives the
annihilated columns, in the elementary abelian setting of
Jennings' theory~\cite{Jennings1941}. A diagonal weighted Gram determinant
and Schwartz's grid bound~\cite[Lemma~1 and Corollary~1]{Schwartz1980},
with related polynomial-identity work of Zippel~\cite{Zippel1979}, supply
enough nondegenerate pairings. Fourier character orbits turn these
projections into actual elements of $K[H]$. The full algebra need not
be semisimple; only its prime-to-$p$ normal abelian subgroup is
diagonalized, in the character-orbit setting of
Clifford~\cite{Clifford1937Announcement}.

\paragraph{From the comparison to scalar elements.}
The proof has three transitions that numerical rank inequalities alone
cannot justify. First, the quotient identifying the cyclic subgroups must
retain the finite vertex groups. Its injectivity proof uses the
$C(4)$--$T(4)$ graphical-complex mechanism described by
Prytu\l a~\cite[Section~8]{Prytula2024}: local short-word restrictions and
triangle-freeness rule out a spherical map with exactly one nontrivial
face. Ordinary graph-of-groups normal form does not by itself establish
this first quotient's injectivity. Section~\ref{sec:gluing} proves it
locally, including the reductions for loops and bridges.

Second, regular ranks must become identities over the same ring. A
further group enlargement $G_1$ has a nonzero central idempotent $f$ in
$R=K[G_1]$. In the ring $S=fR$, whose identity is $f$, the elements
$P_i=fe_i$ admit rectangular matrices
\[
 J_i\in\Mat_{w\times d}(S),\qquad
 Y_i\in\Mat_{d\times w}(S),\qquad
 Y_iJ_i=fI_d,\quad J_iY_i=P_iI_w.
\]
The finite-group blocks in this step have elementary abelian $p$-group
stabilizers, possibly nontrivial and nonnormal. Their group algebras have
nilpotent augmentation ideals, and a finite
change of basis proves exact integer rank identities before coordinates
are paired. Section~\ref{sec:sector} gives the full construction of these
matrices, with their entries in explicitly attached finite subgroups.

Third, a matrix counterexample does not yet give two scalar elements of
the same group algebra. In Section~\ref{op:compression-section}, the
fit and the ordered zero products produce a split injection
\[
 S^{md}\longrightarrow S^{wt}\longrightarrow S^{md},
 \qquad wt<md,
\]
whose composite is the identity. We derive square matrices with a
one-sided inverse and certify the reverse defect by a nonzero kernel
vector. Section~\ref{op:scalar-section} then implements the finite-factor
reduction of Dykema--Juschenko~\cite[Lemma~2.1 and Theorem~2.2]{DykemaJuschenko2015}.
A finite Heisenberg group gives explicit matrix units, their linear
independence makes the encoding injective, and a complementary idempotent
restores the full scalar identity. The shift and modulation model is
classical~\cite[Section~1.2]{GurevichHadani2009}; the composite cyclic
order and finite coefficient field used here are checked directly.

\paragraph{The dynamical consequence.}
A configuration over a finite alphabet $A$ is any function $G\to A$.
A cellular automaton has one finite memory set $E\subset G$ and one
rule applied to $(x(gu))_{u\in E}$ at every $g$. A group is
\emph{surjunctive} when every injective such automaton, over every finite
alphabet, is surjective. Gottschalk proposed this for all groups
\cite{Gottschalk1973}; Gromov and Weiss proved it for sofic groups
\cite{Gromov1999,Weiss2000}. The group-ring connection is recorded by
Elek--Szab\'o and developed in the linear theory of
Ceccherini-Silberstein--Coornaert
\cite{ElekSzabo2004,CSCCoornaert2006}. Phung proved stable finiteness
over every field for surjunctive groups~\cite[Theorem~B]{Phung2023};
subsequent proofs and extensions include~\cite{CSCPhungGroup2025,CSCPhungMonoid2025}.

Section~\ref{op:cellular-section} gives the elementary transfer with its
actual multiplication convention. The coefficients of $b$ define an
injective automaton, those of $a$ give a left inverse, and a nonzero
coefficient of $ba-1$ specifies a missing point indicator. The argument
excludes preimages of arbitrary support. It is a consequence of the
scalar construction, not a further existence theorem.

All matrices below act on columns of free right modules. Matrix entries
therefore multiply coordinates on the left. Transpose is ordinary
transpose over $K$, with no involution. Cardinalities and dimensions in
rank comparisons are integers, not elements of the coefficient field.
For a positive integer $n$, write $[n]=\{1,\ldots,n\}$.
